\documentclass[UTF8,zihao=-4,a4paper,fontset=fandol]{ctexart}
\usepackage{cumcm-paper}

\title{古代玻璃制品的成分分析与鉴别}
\teamnumber{01—074}
\problemnumber{C}
\author{}
\date{}

\begin{document}

\maketitle

\begin{cumcmabstract}
本文围绕古代玻璃制品的成分性数据，建立\textbf{风化校正--类型判别--亚类划分--未知样本鉴别--相关结构比较}的统一流程。对有效成分记录进行闭合归一化和 CLR 变换，先估计风化影响，再将修正后的文物级成分用于分类、聚类和相关分析。\aimark{1}

针对问题一，置换检验和 Firth Logistic 修正表明，\textbf{玻璃类型与表面风化显著相关}：铅钡玻璃原始风化率为 70.0\%，高钾玻璃为 33.3\%；多因素修正后类型和纹饰整体显著。成分上，铅钡玻璃风化后 $SiO_2$ 降低 28.87 个百分点、$PbO$ 升高 22.02 个百分点；高钾玻璃风化后 $SiO_2$ 升高 26.19 个百分点、$K_2O$ 降低 9.02 个百分点。基于类型内 CLR 风化效应还原 32 个风化采样点，49 号、50 号配对验证的 MAE 分别降至 1.848、2.058。

针对问题二，构造铅钡特征相对高钾特征的\textbf{对数比阈值模型}，并用 CLR 最近中心模型复核，两者留一验证准确率均为 98.2\%。高钾玻璃划分为 K-1 高钾钙铝型 15 件和 K-2 高硅含锡特殊型 1 件；铅钡玻璃划分为 LB-1 铅钡均衡型 19 件、LB-2 富钡铜型 6 件和 LB-3 富铅型 15 件，大类分类稳定，亚类需结合风化修正敏感性解释。

针对问题三，8 件未知文物的大类判别在阈值模型、最近中心模型和伪计数扰动下完全一致：\textbf{A1、A6、A7 为高钾 K-1}；A2、A3、A4 为铅钡 LB-3；A5 为铅钡 LB-1；A8 为铅钡 LB-2。其中 A5 的大类稳定，但亚类受风化修正影响，应标注为敏感样本。

针对问题四，高钾与铅钡 CLR 相关矩阵的 Frobenius 差异为 3.832，置换检验 $p=0.0146$，说明两类玻璃\textbf{成分关联结构显著不同}。主要差异来自 $Fe_2O_3$--$CuO$、$CaO$--$CuO$ 和 $Al_2O_3$--$PbO$ 等可靠成分对；含锡、锶、硫等高零值率微量成分只作为特殊样本线索谨慎解释。
\end{cumcmabstract}

\keywords{古代玻璃；成分数据；中心化对数比变换；风化校正；聚类分析}

\CumcmBodyStart

\section{问题重述}

丝绸之路背景下出土的古代玻璃制品同时承载了外来工艺传播和本土材料技术演化的信息。题目给出一批已分类玻璃文物的基本信息、化学成分检测数据以及若干未知类别样本，希望依据化学成分比例和文物外观信息完成风化规律、类别判别、未知样本鉴别和成分关联分析。玻璃成分数据具有两个显著特点：一是各成分为比例数据，理论总和约为 100\%；二是古代玻璃容易受埋藏环境影响发生风化，风化层的成分比例可能明显偏离原始配方。

题目要求分四问完成分析。第一问需要研究表面风化与玻璃类型、纹饰、颜色的关系，比较高钾和铅钡玻璃在风化和无风化状态下的成分变化，并预测风化点的风化前成分。第二问要求依据已分类文物总结高钾与铅钡玻璃的分类规律，并在各大类内部进行亚类划分。第三问要求利用前述规律鉴别表单 3 中未知样本的所属类型，并分析结果敏感性。第四问要求比较不同类别玻璃的成分关联关系及其差异。

本文将表单 1 作为文物级外观与类别信息，将表单 2 作为已分类文物的成分建模数据，将表单 3 作为未知文物预测对象。由于风化会干扰类别判断和相关结构，后续问题均承接第一问得到的风化效应修正结果。

从数据关系看，表单 1 给出的类型标签、纹饰、颜色和表面风化状态属于文物整体属性；表单 2 给出的化学成分属于采样点属性，一件文物可能出现多个采样点，且采样点名称中还包含局部风化状态信息；表单 3 只给出未知类别文物的成分和风化标记。若把三张表简单拼接后直接建模，会出现两个问题：其一，多个采样点会使部分文物在统计上被重复计入；其二，同一风化文物的未风化点与风化点代表不同表面状态，不能仅用文物整体风化标签代替。因此本文在各问中分别明确分析单位：外观因素分析使用文物单位，成分差异分析使用文物--状态单位，类别判别和相关分析使用文物级聚合成分。

题目中的高钾玻璃和铅钡玻璃本质上对应不同助熔体系，前者以钾元素相关成分为主要类别特征，后者以铅、钡元素相关成分为主要类别特征。风化会使部分可溶或易迁移成分发生流失或相对富集，从而掩盖原始配方差异。本文的总体处理思路是先识别并估计风化带来的系统性偏移，再把修正后的成分作为后续分类和相关分析的主要输入。这样做可以使第二问的大类规则、第三问的未知鉴别和第四问的相关结构比较尽量反映文物原始玻璃配方，而不是仅反映表面风化后的检测状态。

\section{问题分析}

本题的核心困难不在于单个算法复杂，而在于成分数据、风化状态和小样本条件同时存在。若直接对原始百分比做普通相关、聚类或分类，成分总和约束会引入虚假负相关；若忽略风化，铅钡玻璃和高钾玻璃的类别差异、未知样本判定和内部亚类都可能被风化层成分扭曲；若把同一文物的多个采样点当作独立样本，则多采样文物会被重复加权。因此，本文先对数据进行有效性筛选和分析单位统一，再根据任务选择文物级或采样点级统计单元。

第一问中，风化影响因素分析以文物为单位，使用列联表、置换检验和 Firth Logistic 回归处理小样本与完全分离问题；成分变化分析以文物--状态单元为单位，按玻璃类型比较风化组和无风化组的成分差异；风化前预测则在 CLR 空间估计类型内风化效应并反向修正。第二问以修正后的文物级成分为输入，先构造可解释的大类判别指标，再在高钾和铅钡内部用 CLR 标准化后的候选成分聚类。第三问直接复用第二问的判别阈值和亚类中心，同时对风化未知样本做类型条件修正验证。第四问在修正后的文物级成分上计算 CLR 相关矩阵，并通过置换检验比较两类玻璃的相关结构差异。

四问之间的信息流为：第一问输出风化前成分预测和类型内风化效应；第二问以该输出为修正数据建立分类规则和亚类中心；第三问利用第二问规则鉴别未知样本；第四问利用问题二的文物级修正成分比较不同类别的成分关联结构。

在模型选择上，本文遵循“先解释规律，再使用复杂方法校验”的原则。对于高钾与铅钡的大类判别，优先构造具有明确化学含义的对数比指标，而不是直接使用黑箱分类器；对于亚类划分，使用层次聚类和 K-means 的一致性、亚类中心成分以及敏感性分析共同判断；对于相关关系，先在 CLR 空间计算相关矩阵，再通过置换检验判断两类结构差异是否超过随机分组可产生的差异。这样既能覆盖题目要求的统计分析，又能保证每个结论都可以回到玻璃成分和文物状态进行解释。

本文还将不确定性分层处理。大类判别由铅、钡、钾三类关键成分主导，样本间差异大，因此作为强结论；风化前预测依赖类型内平均风化效应，可作为后续分析的修正输入，但不解释为单件文物真实原始配方的精确复原；亚类划分和微量成分相关关系受样本量、0 值和特征选择影响更明显，因此在正文中给出敏感性说明。这样的分层有助于避免把所有数值结果都赋予同等确定性。

具体到四问的衔接，第一问承担“校正入口”的作用。如果第一问只给出风化率而不估计风化前成分，那么第二问的聚类可能把“风化程度”误当成“玻璃亚类”，第三问的风化未知样本也可能因表面成分偏移而被错分。第二问承担“规则提取”的作用，它既要给出大类判别的简明规则，也要给出亚类中心，使第三问可以直接外推到未知样本。第三问承担“规则应用”的作用，重点不是重新训练模型，而是检查未知样本在已有规则下是否稳定。第四问承担“结构比较”的作用，它不再只比较各成分均值，而比较不同玻璃体系内部成分之间的组合关系。

因此，本文的技术路线可以理解为一个逐级递进的建模链：先把检测数据整理成适合分析的比例数据；再用风化模型把表面状态影响从原始配方信息中尽量分离；然后用修正后的成分建立分类与亚类；最后用相同的数据基础比较相关结构。所有后续模型都承接同一套预处理和同一套风化修正结果，使四问之间不出现相互矛盾的数据口径。

\section{模型假设}

为使模型与题目数据相匹配，作如下假设：

第一，附件中空白单元格表示该成分未检测到，记为 0。该处理与题意一致，不进行均值填补，以免人为增加微量成分。

第二，风化对成分的影响在同一玻璃类型内部具有可估计的群体平均效应。由于真正的风化前后配对样本很少，本文的风化前预测是类型内反事实修正，而非单件文物的精确物理还原。

第三，经过闭合归一化和 CLR 变换后，不同成分之间的相关、距离和聚类结果可作为比例数据结构的近似刻画。对于大量为 0 的微量成分，相关解释需结合零值率谨慎判断。

第四，置换检验和 Bootstrap 重采样所需的样本可交换性在相应分析层面近似成立。也就是说，在检验某个因素与风化是否有关时，固定样本的类别结构并打乱风化标签可以近似表示“无关联”情形；在评估分类和聚类稳定性时，重复抽样可以近似反映当前样本量下结论的波动范围。

\section{符号说明}

\begin{table}[H]
  \centering
  \caption{主要符号及其含义}
  \label{tab:symbols}
  \begin{tabular}{C{0.18\textwidth}L{0.62\textwidth}C{0.12\textwidth}}
    \toprule
    符号 & 含义 & 单位 \\
    \midrule
    $x_{ij}$ & 第 $i$ 个采样点或文物第 $j$ 种化学成分的原始百分比 & \% \\
    $S_i$ & 第 $i$ 个采样点所有成分百分比之和 & \% \\
    $z_{ij}$ & 闭合归一化后的成分百分比 & \% \\
    $y_{ij}$ & 成分 $z_{ij}$ 经 CLR 变换后的值 & -- \\
    $\varepsilon$ & 处理 0 值的伪计数，主分析取 0.1 & \% \\
    $R_i$ & 铅钡特征相对高钾特征的对数比判别指标 & -- \\
    $r_{ab}^{(t)}$ & 类型 $t$ 内成分 $a,b$ 的 CLR 相关系数 & -- \\
    $D$ & 两类相关矩阵上三角 Frobenius 差异 & -- \\
    \bottomrule
  \end{tabular}
\end{table}

\section{数据预处理与共用成分变换}

\subsection{有效记录筛选与采样点状态识别}

设第 $i$ 条检测记录共有 $p=14$ 种成分，成分总和为
\begin{equation}
  S_i=\sum_{j=1}^{p}x_{ij}.
  \label{eq:sum}
\end{equation}
按照题目要求，仅保留 $85\%\le S_i\le105\%$ 的记录。表单 2 的 69 条采样点记录中有 67 条有效，15 号和 17 号采样点因总和分别为 79.47\%、71.89\% 被剔除。表单 3 的 8 件未知文物成分总和均在有效范围内，全部保留。

采样点名称中含“未风化点”的记录记为无风化，含“严重风化点”的记录记为风化，其余采样点继承表单 1 的文物表面风化状态。该规则允许区分“整件文物表面风化”和“风化文物局部未风化采样点”。

有效性筛选只用于排除明显不满足题目要求的检测记录，不改变保留样本内部各成分的相对结构。15 号和 17 号采样点总和明显偏低，若继续参与分析，会使闭合归一化把所有已检出成分按同一比例放大，进而可能把检测缺失误差误解释为真实成分差异。其余有效记录虽然总和不完全等于 100\%，但处在题目给定容许范围内，经过闭合归一化后可以统一进入比例数据分析。

颜色缺失不进行均值填补或按众数填补，而是作为“未知”类别处理。原因是颜色属于外观分类变量，缺失不一定随机；若用其他颜色填补，会人为改变颜色与风化之间的列联关系。对于后续以成分为主的分类、聚类和相关分析，颜色只在第一问外观因素分析中出现，因此缺失颜色不会影响第二问至第四问的主要成分模型。

\subsection{闭合归一化与 CLR 变换}

由于化学成分是比例数据，本文先对有效记录进行闭合归一化：
\begin{equation}
  z_{ij}=100\frac{x_{ij}}{\sum_{k=1}^{p}x_{ik}}.
  \label{eq:closure}
\end{equation}
对含 0 值的成分，加入伪计数 $\varepsilon$ 后做中心化对数比变换：
\begin{equation}
  y_{ij}=\ln(z_{ij}+\varepsilon)-\frac{1}{p}\sum_{k=1}^{p}\ln(z_{ik}+\varepsilon).
  \label{eq:clr}
\end{equation}
CLR 变换用于距离、聚类、相关和风化效应估计；原始百分比或闭合百分比用于结果解释。

采用 CLR 变换的原因在于，化学成分百分比存在闭合效应：当某一成分比例升高时，即使其他成分的绝对含量没有变化，它们的百分比也会被动下降。如果直接使用原始百分比计算欧氏距离或相关系数，可能得到由总和约束造成的伪相关。CLR 变换把每个成分转化为相对于样本几何均值的对数偏离，更适合用于距离、聚类和相关矩阵。由于题目中存在不少未检出的微量成分，本文加入伪计数 $\varepsilon$，并通过敏感性分析检查该设定对结果的影响。

需要注意的是，CLR 变换后的变量和为 0，单个变量不能像原始百分比一样独立解释。本文在建立模型时使用 CLR 空间保证统计运算合理，在呈现结论时再回到百分比均值、中位数和成分名称进行解释。这样既避免比例数据的技术偏差，又能使结论符合题目中“化学成分含量”的表达习惯。

预处理完成后，本文形成三类核心数据。第一类是采样点级有效数据，用于第一问中识别风化点、未风化点以及预测风化点的风化前成分；第二类是文物--状态级数据，用于比较同一类型内风化状态的成分差异；第三类是文物级聚合数据，用于第二问分类、第三问未知样本判别和第四问相关结构比较。三类数据的区别在于统计单位不同，但均来自同一套有效性筛选和闭合归一化规则。

这种分层数据结构使模型能够兼顾题目给出的不同信息粒度。表面风化与纹饰、颜色的关系必须在文物层面讨论，因为纹饰和颜色不是采样点属性；风化成分差异必须保留状态信息，因为同一文物可能同时包含风化点和未风化点；大类分类和亚类划分则应回到文物级，因为题目要求鉴别的是文物类型而不是某个检测点类型。若不区分这些单位，统计显著性和分类准确率都可能被重复采样放大。

\subsection{全题共用数据流水线}

为使四问的计算口径一致，本文在正式建模前把原始表格转换成四个中间对象。文物属性表来自表单 1，以文物编号为键，记录玻璃类型、纹饰、颜色和表面风化状态；已知类别采样点成分表来自表单 2，先按成分总和筛选有效记录，再识别采样点风化状态，并保存闭合百分比矩阵和 CLR 矩阵；文物--状态矩阵按“文物编号 + 风化状态”分组求均值，用于比较风化与无风化状态；未知样本矩阵来自表单 3，只在第三问作为待鉴别对象进入固定模型。

第一问完成后，还会生成一份文物级风化修正矩阵。该矩阵把风化采样点替换为风化前预测成分，无风化采样点保留原始有效成分，再按文物编号等权聚合。第二问的大类分类和亚类划分、第三问的未知样本鉴别、第四问的相关结构比较都读取这份矩阵，因此后续问题不会各自使用不同版本的成分数据。

这种流水线还规定了各问之间的输入输出关系。第一问的输出不是一句风化结论，而是风化效应向量和风化前成分预测表；第二问读取该预测表，输出大类阈值、最近中心和亚类中心；第三问读取第二问的固定规则，输出 A1--A8 的类型、亚类和敏感性标记；第四问读取第二问的文物级修正矩阵，输出两类相关矩阵、整体差异检验和差异成分对。这样，正文中的每个结果都能追溯到前一阶段的明确输出。

所有中间表均保留统一编号字段。已知样本使用文物编号连接表单 1 与表单 2，未知样本使用 A1--A8 连接表单 3 与第三问结果表；成分列顺序在预处理、分类、聚类和相关分析中保持一致。这样可以避免在多次变换后出现样本错位或成分错列，保证每个数值结果都能回到原始附件中的对应记录。

\section{问题一：风化影响、成分规律与风化前预测}

\subsection{任务拆分与数据口径}

第一问先把题目拆成三个连续的小问题：风化状态受哪些外观因素影响、风化前后成分怎样改变、已风化采样点应怎样反推风化前成分。三者不能并列求解，因为风化因素分析给出风化现象的统计背景，成分差异比较给出修正方向，风化前预测才是后续分类和相关分析真正需要的输入。

文物级风化因素分析只使用表单 1 中每件文物的一条记录，因变量为是否表面风化，自变量为玻璃类型、纹饰和颜色；若颜色缺失，则作为“未知”类别进入列联表和回归模型，不进行颜色填补。成分差异分析则转到表单 2 的采样点层面：含“未风化点”者记为无风化，含“严重风化点”者记为风化，其余采样点继承文物表面风化状态。这样，风化因素判断和成分修正各自使用与任务相匹配的统计单位。

单因素置换检验的执行方式为：固定某一因素的类别计数和风化样本总数，将风化标签随机打乱 5000 次，每次重新计算列联表卡方统计量。若观测统计量为 $\chi^2_{\mathrm{obs}}$，第 $b$ 次置换统计量为 $\chi^2_b$，则置换 $p$ 值取
\begin{equation}
  p_{\mathrm{perm}}=\frac{1+\sum_{b=1}^{B}I(\chi^2_b\ge \chi^2_{\mathrm{obs}})}{B+1}.
  \label{eq:perm-p}
\end{equation}
这样得到的 $p$ 值不依赖大样本卡方近似，更适合本题部分类别频数较小的情况。

成分比较部分先把同一文物、同一风化状态下的多个采样点聚合为一个均值向量。若第 $m$ 件文物在状态 $s$ 下有 $q_m$ 个有效采样点，则其第 $j$ 种成分的状态均值为
\begin{equation}
  \bar x_{msj}=\frac{1}{q_m}\sum_{\ell=1}^{q_m}x_{\ell j}.
  \label{eq:state-mean}
\end{equation}
随后在同一玻璃类型内部比较风化状态 $W$ 与无风化状态 $U$。这一口径保证一件文物在同一状态下只贡献一个均值向量，避免多采样点文物在检验中过度加权。

风化前预测输出分为采样点级和文物级两层。采样点级输出保留每个风化采样点的 14 种预测成分，用于检查修正方向；文物级输出则把同一文物的有效点等权平均，用于第二问分类和第四问相关分析。也就是说，第一问结束时不仅给出“哪些因素与风化相关”的统计结论，还实际生成了后续模型需要使用的修正成分矩阵。

第一问的三类输出按同一编号体系保存。风化影响因素表记录每个类型、纹饰和颜色类别的样本数、风化数、风化率、置换 $p$ 值和 Cramer's $V$；成分变化表记录每个玻璃类型内风化组均值、无风化组均值、均值差、置换 $p$ 值和 BH 校正后的 $q$ 值；风化前预测表记录每个风化采样点的原始成分、CLR 修正量和逆变换后的预测成分。正文后续使用这些输出时，不再重新从原始采样点直接计算，以保证问题一的修正结果能连续传递到问题二至问题四。

实际落表时，本文先生成文物级风化因素表，再生成采样点状态表和文物--状态均值表，最后生成风化前预测表。这个顺序对应第一问的求解逻辑：先判断风化现象是否与类型和外观有关，再在同一类型内部估计风化造成的成分位移，最后把该位移反向用于风化点预测。若跳过中间的文物--状态均值表，风化效应就会被多采样点文物放大；若直接用全部样本估计统一修正量，高钾和铅钡两类玻璃的不同风化方向会被混在一起。

\subsection{表面风化影响因素分析}

第一问首先以 58 件文物为单位建立类型、纹饰、颜色与表面风化状态的列联表。对第 $k$ 个类别定义风化率
\begin{equation}
  R_k=\frac{n_{k,W}}{n_k},
  \label{eq:weather-rate}
\end{equation}
其中 $n_{k,W}$ 为该类别风化文物数，$n_k$ 为该类别文物总数。为避免小样本颜色类别下渐近卡方检验不稳定，本文使用固定类别、随机打乱风化标签的置换检验，并同时计算 Cramer's $V$ 作为关联强度。

\begin{table}[H]
  \centering
  \caption{表面风化与类型、纹饰、颜色的单因素关联检验}
  \label{tab:q1-association}
  \begin{tabular}{lrrrr}
    \toprule
    因素 & 卡方统计量 & 自由度 & 置换 $p$ 值 & Cramer's $V$ \\
    \midrule
    玻璃类型 & 6.880 & 1 & 0.0125 & 0.344 \\
    纹饰 & 4.957 & 2 & 0.1034 & 0.292 \\
    颜色 & 9.432 & 8 & 0.3172 & 0.403 \\
    \bottomrule
  \end{tabular}
\end{table}

表 \ref{tab:q1-association} 给出第一步结果：未控制其他因素时，玻璃类型与风化显著相关。铅钡玻璃 40 件中 28 件风化，风化率为 70.0\%；高钾玻璃 18 件中 6 件风化，风化率为 33.3\%。纹饰和颜色在单因素下未达到 0.05 显著性水平，但纹饰 B 的 6 件样本全部风化，因此下一步需要把类型、纹饰和颜色放入同一模型共同修正。

考虑纹饰 B、黑色等类别存在完全或近完全分离，本文采用 Firth 偏差修正 Logistic 回归：
\begin{equation}
  \operatorname{logit}(p_i)=\beta_0+\beta_TT_i+\beta_BD_{Bi}+\beta_CD_{Ci}+\sum_k\gamma_kC_{ki},
  \label{eq:firth-model}
\end{equation}
其中 $p_i$ 为第 $i$ 件文物表面风化概率，$T_i$ 表示类型，$D_{Bi},D_{Ci}$ 表示纹饰虚拟变量，$C_{ki}$ 表示颜色虚拟变量。Firth 模型最大化惩罚似然
\begin{equation}
  \ell_F(\bm\beta)=\ell(\bm\beta)+\frac{1}{2}\ln|\bm I(\bm\beta)|,
  \label{eq:firth}
\end{equation}
从而在小样本分离情形下得到有限估计。

\begin{table}[H]
  \centering
  \caption{多因素修正后风化影响因素整体检验与调整概率}
  \label{tab:q1-firth}
  \begin{tabular}{lrrl}
    \toprule
    因素或类别 & Bootstrap $p$ 值 & 调整后风化概率 & 95\% 区间 \\
    \midrule
    类型整体 & 0.0002 & -- & -- \\
    纹饰整体 & 0.0006 & -- & -- \\
    颜色整体 & 0.0574 & -- & -- \\
    铅钡 & -- & 73.6\% & 59.9\%--80.7\% \\
    高钾 & -- & 13.7\% & 11.0\%--28.2\% \\
    纹饰 A & -- & 34.0\% & 22.5\%--53.0\% \\
    纹饰 B & -- & 95.0\% & 74.8\%--97.9\% \\
    纹饰 C & -- & 58.7\% & 46.3\%--73.3\% \\
    \bottomrule
  \end{tabular}
\end{table}

Firth Logistic 的计算过程为：先用式 \eqref{eq:firth-model} 设定风化概率，再用式 \eqref{eq:firth} 的惩罚似然估计参数，随后用 Bootstrap 对类型、纹饰、颜色三个因素分别做整体检验，最后在平均意义上计算各类别的调整后风化概率。表 \ref{tab:q1-firth} 显示，控制其他变量后，类型整体 $p=0.0002$，纹饰整体 $p=0.0006$，均达到显著水平；铅钡调整后风化概率为 73.6\%，高钾为 13.7\%，纹饰 B 调整后风化概率为 95.0\%。因此本部分输出为：类型和纹饰与表面风化显著相关，颜色只作为辅助因素。

\subsection{分类型风化成分差异}

成分差异分析以文物--状态单元为样本。对于同一文物在同一风化状态下的多个有效采样点，先对 14 种成分取均值，避免多采样文物重复加权。最终得到铅钡风化 21 个、铅钡无风化 21 个、高钾风化 6 个、高钾无风化 10 个文物--状态单元。

对玻璃类型 $t$ 和成分 $j$ 定义风化均值差
\begin{equation}
  \Delta_{tj}=\bar x_{tj}^{(W)}-\bar x_{tj}^{(U)}.
  \label{eq:component-diff}
\end{equation}
在每个类型内部随机打乱风化标签进行置换检验，并用 Benjamini--Hochberg 方法控制 14 个成分的多重比较。

\begin{table}[H]
  \centering
  \caption{风化与无风化状态下显著变化的主要成分}
  \label{tab:q1-component}
  \begin{tabular}{llrrrr}
    \toprule
    类型 & 成分 & 风化均值 & 无风化均值 & 差值 & $q$ 值 \\
    \midrule
    铅钡 & $SiO_2$ & 26.85 & 55.73 & -28.87 & 0.0003 \\
    铅钡 & $PbO$ & 43.37 & 21.36 & 22.02 & 0.0003 \\
    铅钡 & $P_2O_5$ & 4.91 & 1.11 & 3.79 & 0.0007 \\
    铅钡 & $CaO$ & 2.56 & 1.22 & 1.34 & 0.0082 \\
    铅钡 & $Na_2O$ & 0.27 & 1.57 & -1.30 & 0.0305 \\
    高钾 & $SiO_2$ & 93.96 & 67.77 & 26.19 & 0.0035 \\
    高钾 & $K_2O$ & 0.54 & 9.57 & -9.02 & 0.0040 \\
    高钾 & $CaO$ & 0.87 & 5.73 & -4.86 & 0.0066 \\
    高钾 & $Al_2O_3$ & 1.93 & 6.41 & -4.48 & 0.0040 \\
    高钾 & $MgO$ & 0.20 & 1.05 & -0.86 & 0.0258 \\
    高钾 & $Fe_2O_3$ & 0.27 & 1.79 & -1.52 & 0.0258 \\
    \bottomrule
  \end{tabular}
\end{table}

成分差异的计算顺序为：先按“文物--风化状态”聚合采样点，再在每个玻璃类型内分别计算式 \eqref{eq:component-diff} 的风化均值差，随后在类型内打乱风化标签做置换检验，最后对 14 个成分的 $p$ 值进行 BH 校正。表 \ref{tab:q1-component} 是该流程的直接输出。铅钡玻璃风化后 $SiO_2$ 降低 28.87 个百分点，$PbO$ 升高 22.02 个百分点；高钾玻璃风化后 $SiO_2$ 升高 26.19 个百分点，$K_2O$ 降低 9.02 个百分点。两类方向相反，因此后续不能使用一个全样本风化修正向量，而必须在高钾和铅钡内部各自估计风化效应。

\subsection{风化前成分预测}

为了预测风化点的风化前成分，本文在每种玻璃类型内部估计风化组与无风化组的中位 CLR 差：
\begin{equation}
  d_{tj}=\operatorname{median}(y_{ij}\mid t,W)-\operatorname{median}(y_{ij}\mid t,U).
  \label{eq:weather-effect}
\end{equation}
对类型为 $t_i$ 的风化点，其风化前 CLR 向量估计为
\begin{equation}
  \hat y_{ij}^{(U)}=y_{ij}^{(W)}-d_{t_i j}.
  \label{eq:preweather-clr}
\end{equation}
最后用逆 CLR 变换闭合到 100\%：
\begin{equation}
  \hat x_{ij}^{(U)}=100\frac{\exp(\hat y_{ij}^{(U)})}{\sum_{k=1}^{p}\exp(\hat y_{ik}^{(U)})}.
  \label{eq:inverse-clr}
\end{equation}
主分析取 $\varepsilon=0.1$，并将单个成分的 CLR 修正幅度限制在 $[-2.5,2.5]$ 内；实际没有成分触发该限制。该方法对 32 个有效风化采样点完成风化前预测。

风化前预测的实际计算为：先将所有有效成分闭合归一化并做 CLR 变换；再按玻璃类型分别用式 \eqref{eq:weather-effect} 计算风化组与无风化组的中位 CLR 差；对每个风化点用式 \eqref{eq:preweather-clr} 扣除对应类型的风化效应；最后用式 \eqref{eq:inverse-clr} 转回百分比成分并闭合到 100\%。该流程保证了每个风化点都有一组可直接进入第二问分类模型的风化前成分。

\begin{table}[H]
  \centering
  \caption{风化前预测的敏感性与配对样本验证}
  \label{tab:q1-validation}
  \begin{tabular}{lrrr}
    \toprule
    检验项目 & 指标一 & 指标二 & 指标三 \\
    \midrule
    伪计数 0.05 相对 0.10 & 平均变化 0.201 & 最大变化 4.310 & 95\% 分位 0.888 \\
    伪计数 0.50 相对 0.10 & 平均变化 0.829 & 最大变化 5.897 & 95\% 分位 3.677 \\
    49 号铅钡配对样本 & 原始 MAE 3.696 & 修正后 MAE 1.848 & 改善 1.848 \\
    50 号铅钡配对样本 & 原始 MAE 3.716 & 修正后 MAE 2.058 & 改善 1.659 \\
    \bottomrule
  \end{tabular}
\end{table}

表 \ref{tab:q1-validation} 用两个检查验证预测过程。第一，改变伪计数后重新计算完整预测流程，伪计数 0.05 相对 0.10 的平均变化为 0.201，说明主设定附近预测稳定；第二，用 49、50 号铅钡配对样本比较修正前后 MAE，修正后分别降至 1.848 和 2.058。第一问最终输出为两部分：一是铅钡玻璃和纹饰 B 更易风化，二是 32 个风化采样点的风化前成分预测。第二问以后只对风化点替换为预测成分，无风化点保留原始有效成分，再按文物等权聚合。

\section{问题二：分类规律与亚类划分}

\subsection{文物级训练矩阵构建}

第二问从第一问的风化前预测结果接入，不再直接使用混有风化影响的原始采样点。设第 $i$ 件文物的文物级闭合成分向量为 $\bm x_i=(x_{i1},\ldots,x_{ip})$，已知大类标签为 $g_i\in\{\text{高钾},\text{铅钡}\}$。对无风化采样点保留原始有效成分，对风化采样点替换为第一问反推的风化前成分；同一文物若有多个有效采样点，则先在百分比空间等权平均，再重新闭合归一化。由此得到 56 件文物、14 种成分的训练矩阵，其中每一行对应一个已知类别文物。

\subsection{大类阈值模型建立}

大类判别先抓住高钾和铅钡最稳定的成分差异。高钾玻璃以 $K_2O$ 为主要助熔相关成分，铅钡玻璃以 $PbO$、$BaO$ 为主要特征成分，因此定义铅钡特征相对高钾特征的对数比指标 $R_i$。当 $R_i$ 越大，样本越接近铅钡；当 $R_i$ 越小，样本越接近高钾。
\begin{equation}
  R_i=\ln\frac{PbO_i+BaO_i+\varepsilon}{K_2O_i+\varepsilon}.
  \label{eq:major-ratio}
\end{equation}
若 $R_i$ 大于阈值，则判为铅钡；否则判为高钾。

大类阈值搜索按以下计算表执行。首先对每件文物计算 $R_i$；其次按 $R_i$ 升序排列样本，取相邻两个不同 $R_i$ 的中点作为候选阈值 $\tau$；再次对每个 $\tau$ 计算
\begin{equation}
  \mathrm{Se}(\tau)=\frac{TP}{TP+FN},\quad
  \mathrm{Sp}(\tau)=\frac{TN}{TN+FP},\quad
  J(\tau)=\mathrm{Se}(\tau)+\mathrm{Sp}(\tau)-1.
  \label{eq:youden}
\end{equation}
其中铅钡作为阳性类。最终选择准确率最高且 $J(\tau)$ 最大的候选阈值，得到 $\tau=0.501$。留一验证时，第 $i$ 件文物先从训练集中移出，用其余 55 件重新完成阈值搜索，再判定第 $i$ 件文物，循环 56 次后汇总准确率。

\subsection{CLR 最近中心复核}

CLR 最近中心复核模型的计算也固定为留一形式。每次留出一件文物后，只在训练部分计算高钾中心 $\bm c_H$ 和铅钡中心 $\bm c_L$。对留出样本的 CLR 标准化向量 $\bm y_i$，计算
\begin{equation}
  d_H(i)=\|\bm y_i-\bm c_H\|_2,\qquad
  d_L(i)=\|\bm y_i-\bm c_L\|_2.
  \label{eq:nearest-center}
\end{equation}
若 $d_H(i)<d_L(i)$，判为高钾；否则判为铅钡。该模型不替代对数比阈值，而是检查多成分整体距离是否支持同一大类判断。

\subsection{类内亚类聚类}

亚类划分的具体运算先在每个大类内部独立完成。对候选成分矩阵做 CLR 变换后，再对每一列进行均值为 0、标准差为 1 的标准化，得到聚类矩阵 $\bm Z$。Ward 层次聚类使用欧氏距离和最小类内平方和增加准则，给出合并树；K-means 在候选类数 $k$ 下重复初始化，取类内平方和最小的结果。亚类数不是凭文字命名确定，而是综合轮廓系数、类内平方和拐点、层次树分裂位置和类中心成分差异确定。

亚类稳定性用调整 Rand 指数衡量。若两种设置下得到的划分分别为 $P$ 和 $Q$，ARI 接近 1 表示两个划分几乎一致，接近 0 表示一致性接近随机。本文分别比较“原始有效数据 vs. 风化修正数据”“伪计数 0.05 或 0.50 vs. 主设定 0.10”“删除单个候选成分 vs. 全特征”。这些计算使第二问不仅给出分类标签，还说明哪些标签可直接外推到第三问，哪些只应作为敏感亚类解释。

类内聚类完成后，本文并不直接把聚类编号作为亚类名称。先计算每个聚类中心在原始成分百分比空间中的中位水平，再与本大类总体中位水平比较，找出相对偏高或偏低的关键成分。高钾类中，主体类的 $K_2O$、$CaO$、$Al_2O_3$ 较稳定，因此命名为 K-1 高钾钙铝型；单样本类 18 号表现出较高 $SiO_2$ 和 $SnO_2$，只命名为 K-2 高硅含锡特殊型。铅钡类中，LB-1 的 $PbO$、$BaO$ 较均衡，LB-2 的 $BaO$、$CuO$ 同时偏高，LB-3 的 $PbO$ 偏高而 $BaO$ 较低。这样，亚类名称由中心成分计算得到，而不是先设定名称再解释样本。

第二问向第三问输出三组固定对象。第一组是大类阈值 $\tau=0.501$ 和式 \eqref{eq:major-ratio}，用于未知样本的初步大类判断；第二组是高钾、铅钡两个 CLR 最近中心及其所用成分列，用于复核阈值判断；第三组是 K-1、K-2、LB-1、LB-2、LB-3 五个亚类中心，用于第三问的亚类距离比较。第三问只读取这些对象，不重新选择成分、不重新搜索阈值，也不改变亚类数。

\subsection{大类判别结果}

在 56 件有效文物中，高钾 16 件、铅钡 40 件。按式 \eqref{eq:major-ratio} 搜索得到主阈值 0.501，并用 CLR 最近中心模型复核，两种模型的验证结果见表 \ref{tab:q2-major}。

\begin{table}[H]
  \centering
  \caption{高钾与铅钡大类分类模型表现}
  \label{tab:q2-major}
  \begin{tabular}{lrrrr}
    \toprule
    模型 & 准确率 & 铅钡灵敏度 & 高钾特异度 & 错分数 \\
    \midrule
    对数比阈值模型（样本内） & 98.2\% & 100.0\% & 93.8\% & 1 \\
    对数比阈值模型（留一） & 98.2\% & 100.0\% & 93.8\% & 1 \\
    CLR 最近中心模型（留一） & 98.2\% & 97.5\% & 100.0\% & 1 \\
    \bottomrule
  \end{tabular}
\end{table}

表 \ref{tab:q2-major} 是大类模型的验证输出。对数比阈值模型和 CLR 最近中心模型留一准确率均为 98.2\%，错分数均为 1。因此大类规律写为：铅钡玻璃以 $PbO$、$BaO$ 高而 $K_2O$ 低为特征，高钾玻璃以 $K_2O$ 较高且 $PbO$、$BaO$ 接近未检出为特征。第三问将直接复用阈值 0.501 和最近中心模型。

这个结果表的生成顺序为：先在全部 56 件文物上得到样本内判别结果，再在留一循环中记录每次被留出样本的预测标签，最后把阈值模型和最近中心模型的灵敏度、特异度、准确率和错分数合并。留一验证中，留出样本不参与阈值搜索和中心计算，因此表 \ref{tab:q2-major} 的 98.2\% 不是简单的样本内拟合，而是对已知样本进行逐件外推后的结果。

\subsection{亚类划分方法}

亚类划分只在已知大类内部进行。高钾候选成分取 $SiO_2$、$K_2O$、$CaO$、$MgO$、$Al_2O_3$、$Fe_2O_3$ 和 $CuO$；铅钡候选成分取 $SiO_2$、$PbO$、$BaO$、$P_2O_5$、$SrO$ 和 $CuO$。计算时先对候选成分做 CLR 变换和标准化，再分别计算 Ward 层次聚类树，并在候选类数下运行 K-means。根据轮廓系数、类内平方和、树状结构和类中心成分，最终取高钾 2 类、铅钡 3 类。亚类命名只依据类中心成分，不额外推断产地或年代。

\begin{table}[H]
  \centering
  \caption{高钾与铅钡玻璃亚类划分结果}
  \label{tab:q2-subclass}
  \begin{tabular}{llrL{0.43\textwidth}}
    \toprule
    大类 & 亚类 & 件数 & 主要文物编号与成分特征 \\
    \midrule
    高钾 & K-1 高钾钙铝型 & 15 & 以 01、03、04 等 15 件为主体；中位 $K_2O=10.55$、$CaO=7.41$、$Al_2O_3=6.28$。 \\
    高钾 & K-2 高硅含锡特殊型 & 1 & 18；$SiO_2=81.71$、$SnO_2=2.43$，样本数为 1，作为特殊样本而非稳定普遍亚类解释。 \\
    铅钡 & LB-1 铅钡均衡型 & 19 & $PbO=19.66$、$BaO=6.85$，$SiO_2$ 相对较高，代表较典型的铅钡结构。 \\
    铅钡 & LB-2 富钡铜型 & 6 & 08,11,20,24,26,37；$BaO=26.58$、$CuO=4.65$，钡和铜显著偏高。 \\
    铅钡 & LB-3 富铅型 & 15 & $PbO=22.71$、$BaO=4.38$，铅含量相对高而钡含量较低。 \\
    \bottomrule
  \end{tabular}
\end{table}

表 \ref{tab:q2-subclass} 是亚类划分输出。高钾分为 K-1 高钾钙铝型 15 件和 K-2 高硅含锡特殊型 1 件；铅钡分为 LB-1 铅钡均衡型 19 件、LB-2 富钡铜型 6 件和 LB-3 富铅型 15 件。命名过程为：先读取各类中心和中位成分，再用最突出的成分特征形成中文名称，例如 LB-2 同时具有较高 $BaO$ 和 $CuO$，因此命名为富钡铜型。

\subsection{合理性与敏感性}

敏感性检验按三个操作执行：第一，用原始有效数据和风化修正数据分别聚类，比较调整 Rand 指数；第二，将伪计数从 0.10 改为 0.05 和 0.50，重复完整聚类流程；第三，每次删除一个候选成分后重新聚类，判断亚类是否由单一变量支配。结果显示，大类判别是强结论；高钾亚类在原始与修正数据下 ARI 为 1.000，较稳定；铅钡亚类对应 ARI 为 0.434，需在第三问中对边界样本标注敏感性。第二问向第三问输出三项内容：大类阈值 0.501、CLR 最近中心模型和五个亚类的标准化中心。

第二问最终表的生成按固定顺序完成。先根据大类标签把 56 件文物拆成高钾组和铅钡组，再在各组内部读取聚类标签；随后计算每个亚类的中心成分、样本数和相对本大类均值的偏离量，形成亚类名称；最后把每件文物的文物编号、大类、亚类、到本亚类中心距离和敏感性标记合并为输出表。这样表中的每个亚类标签都能回溯到一个中心向量和一组距离，而不是仅由人工描述决定。

\section{问题三：未知类别玻璃文物鉴别}

\subsection{未知样本有效性与大类初判}

第三问只把第二问已经确定的规则用于表单 3，不把未知样本加入训练集重新调阈值或重定亚类。对每件未知文物先计算成分总和 $S_i$，总和满足有效范围后进行闭合归一化和 CLR 变换。随后代入式 \eqref{eq:major-ratio} 得到 $R_i$，与第二问阈值 0.501 比较，形成第一组大类判断。

为避免单一阈值误判，每个未知样本还进入 CLR 最近中心模型。若其到高钾中心的距离 $d_H(i)$ 小于到铅钡中心的距离 $d_L(i)$，则复核结果为高钾；反之为铅钡。阈值判断和最近中心判断一致时，进入后续亚类判定；若二者不一致，则把样本列为低置信度样本。

\subsection{风化样本的类型条件修正}

表单 3 中 A2、A5、A6、A7 标记为风化样本，不能只按原始风化成分判定。本文对这四件文物分别计算两条修正路径：一条假设其为高钾，扣除高钾风化效应向量后重新闭合；另一条假设其为铅钡，扣除铅钡风化效应向量后重新闭合。每条路径都重新计算 $R_i$、$d_H(i)$ 和 $d_L(i)$。

若原始成分与类型条件修正后的大类结果相同，则认为大类对风化修正稳定；若大类相同但亚类中心距离排序改变，则保留修正后亚类并标记为敏感。A5 就属于后一种情况：大类始终为铅钡，但亚类距离受风化修正影响。

\subsection{亚类距离判定}

亚类判别使用限制搜索范围的原则。若大类为高钾，则只计算到 K-1、K-2 两个中心的距离：
\begin{equation}
  h(i)=\arg\min_{k\in\{K\text{-}1,K\text{-}2\}}\|\bm z_i-\bm \mu_k\|_2.
  \label{eq:subclass-k}
\end{equation}
若大类为铅钡，则只计算到 LB-1、LB-2、LB-3 的距离：
\begin{equation}
  l(i)=\arg\min_{k\in\{LB\text{-}1,LB\text{-}2,LB\text{-}3\}}\|\bm z_i-\bm \mu_k\|_2.
  \label{eq:subclass-lb}
\end{equation}
其中 $\bm z_i$ 为未知样本在相应候选成分上的 CLR 标准化向量，$\bm \mu_k$ 为第二问得到的亚类中心。这样可以避免跨大类中心造成的误配。

\subsection{逐样本判别与稳定性标记}

无风化样本 A1、A3、A4、A8 直接使用原始闭合成分判别。A1 的铅钡特征弱、与高钾中心距离更近，归入高钾 K-1；A3、A4 到 LB-3 中心距离最小，归入铅钡富铅型；A8 与 LB-2 中心距离最小，归入富钡铜型。风化样本 A2、A5、A6、A7 使用类型条件修正后的成分给出亚类，其中 A2 归入 LB-3，A6、A7 归入 K-1，A5 归入 LB-1 但标记为亚类敏感。

稳定性标记由三类检查合并得到：阈值模型与最近中心模型是否一致，风化条件修正前后大类是否一致，伪计数取 0.05、0.10、0.50 时大类是否一致。若三类检查均不改变大类，则报告“大类稳定”；若仅亚类发生变化，则只在亚类层面标注敏感。

\subsection{未知文物鉴别结果}

\begin{table}[H]
  \centering
  \caption{表单 3 未知文物类型与亚类鉴别结果}
  \label{tab:q3-final}
  \begin{tabular}{ccccc}
    \toprule
    文物 & 表面风化 & 判别类型 & 亚类 & 稳定性说明 \\
    \midrule
    A1 & 无风化 & 高钾 & K-1 高钾钙铝型 & 大类与亚类稳定 \\
    A2 & 风化 & 铅钡 & LB-3 富铅型 & 大类与亚类稳定 \\
    A3 & 无风化 & 铅钡 & LB-3 富铅型 & 大类与亚类稳定 \\
    A4 & 无风化 & 铅钡 & LB-3 富铅型 & 大类与亚类稳定 \\
    A5 & 风化 & 铅钡 & LB-1 铅钡均衡型 & 大类稳定，亚类受风化修正影响 \\
    A6 & 风化 & 高钾 & K-1 高钾钙铝型 & 大类与亚类稳定 \\
    A7 & 风化 & 高钾 & K-1 高钾钙铝型 & 大类与亚类稳定 \\
    A8 & 无风化 & 铅钡 & LB-2 富钡铜型 & 大类与亚类稳定 \\
    \bottomrule
  \end{tabular}
\end{table}

表 \ref{tab:q3-final} 是逐样本流程的最终输出。A1、A6、A7 判为高钾 K-1；A2、A3、A4 判为铅钡 LB-3；A5 判为铅钡 LB-1；A8 判为铅钡 LB-2。其中 A5 的大类在阈值模型、最近中心模型和风化修正下均为铅钡，但亚类由原始成分下的富铅型转为风化修正后的铅钡均衡型，因此标注为“亚类受风化修正影响”。第三问最终结论按置信层级报告：大类全部稳定；除 A5 外亚类稳定；A5 大类稳定但亚类敏感。

\section{问题四：不同类别的成分关联关系比较}

\subsection{CLR 相关矩阵构建}

第四问比较两类玻璃内部成分关系，不再把类别作为预测目标。输入矩阵沿用第二问的文物级风化修正成分，按已知标签拆成高钾组和铅钡组；每组都对 14 种成分做 CLR 变换，再计算同组内部的成分相关矩阵。

设类型 $t$ 的 CLR 矩阵为 $\bm Y^{(t)}$，其第 $a$、$b$ 列分别表示两个成分在该类文物中的 CLR 值。对类型 $t$ 内成分 $a,b$，相关系数定义为
\begin{equation}
  r_{ab}^{(t)}=
  \frac{\sum_i(y_{ia}-\bar y_a)(y_{ib}-\bar y_b)}
  {\sqrt{\sum_i(y_{ia}-\bar y_a)^2}\sqrt{\sum_i(y_{ib}-\bar y_b)^2}}.
  \label{eq:corr}
\end{equation}
对所有 $a<b$ 计算一次，共得到 91 个成分对相关系数。若 $r_{ab}^{(t)}>0$，说明两成分在该类玻璃中相对含量同步变化；若 $r_{ab}^{(t)}<0$，说明二者呈替代或反向变化。

由于微量成分可能大量为 0，每个成分对同时记录最大零值率
\begin{equation}
  z_{ab}=\max(z_a,z_b),
  \label{eq:zero-rate}
\end{equation}
其中 $z_a$、$z_b$ 分别为两个成分在该类型内的原始 0 值比例。正文优先解释 $z_{ab}$ 较低且相关系数绝对值较大的成分对。

\subsection{整体相关结构差异检验}

整体结构差异只使用相关矩阵上三角元素，避免重复计算对称位置。具体操作是将高钾相关矩阵和铅钡相关矩阵的上三角拉直为两个长度为 91 的向量，再计算式 \eqref{eq:frobenius} 的欧氏距离。置换检验时，每次随机重排“高钾/铅钡”标签，但保持 16 件和 40 件的组容量不变；重排后重新计算两组相关矩阵和 $D_b$。最终整体 $p$ 值按式 \eqref{eq:perm-p} 的同样形式计算。

\subsection{差异成分对筛选}

单个成分对差异检验在整体显著后进行。对成分对 $(a,b)$，先计算观测差异
\begin{equation}
  \delta_{ab}=r_{ab}^{(H)}-r_{ab}^{(L)}.
  \label{eq:corr-diff}
\end{equation}
再在每次置换中计算 $\delta_{ab,b}$，用 $|\delta_{ab,b}|\ge|\delta_{ab}|$ 的比例得到双侧置换 $p$ 值。由于 91 个成分对同时检验，最后用 BH 方法控制多重比较，只在正文列出差异较大、零值率较低且解释稳定的成分对。

\subsection{类型内部强相关关系}

\begin{table}[H]
  \centering
  \caption{不同类型玻璃内部的主要可靠强相关成分对}
  \label{tab:q4-top}
  \begin{tabular}{llrr}
    \toprule
    类型 & 成分对 & 相关系数 & 最大零值率 \\
    \midrule
    高钾 & $Fe_2O_3$--$CuO$ & 0.781 & 0.063 \\
    高钾 & $CaO$--$SrO$ & -0.701 & 0.313 \\
    高钾 & $SiO_2$--$Fe_2O_3$ & -0.665 & 0.063 \\
    高钾 & $CaO$--$CuO$ & 0.622 & 0.063 \\
    铅钡 & $Na_2O$--$P_2O_5$ & -0.704 & 0.250 \\
    铅钡 & $MgO$--$BaO$ & -0.526 & 0.200 \\
    铅钡 & $CuO$--$BaO$ & 0.482 & 0.025 \\
    铅钡 & $Al_2O_3$--$BaO$ & -0.475 & 0.000 \\
    铅钡 & $Al_2O_3$--$PbO$ & -0.441 & 0.000 \\
    \bottomrule
  \end{tabular}
\end{table}

表 \ref{tab:q4-top} 是类内相关矩阵筛选结果。高钾玻璃中 $Fe_2O_3$--$CuO$ 相关系数为 0.781，且最大零值率仅 0.063，是较可靠的正相关；$SiO_2$--$Fe_2O_3$ 为 -0.665，说明硅质基体和铁相关杂质呈反向变化。铅钡玻璃中 $CuO$--$BaO$ 为 0.482，最大零值率 0.025，与第二问 LB-2 富钡铜型相互对应；$Al_2O_3$--$PbO$ 和 $Al_2O_3$--$BaO$ 均为负相关且零值率为 0，说明铝硅矿物杂质与铅钡助熔剂呈替代式比例关系。

\subsection{两类相关结构差异检验}

设高钾相关矩阵为 $\bm R_H$，铅钡相关矩阵为 $\bm R_L$，用上三角 Frobenius 距离衡量整体差异：
\begin{equation}
  D=\left[\sum_{a<b}(r_{ab}^{(H)}-r_{ab}^{(L)})^2\right]^{1/2}.
  \label{eq:frobenius}
\end{equation}
保持两类样本量不变，随机置换类型标签 5000 次，得到 $D$ 的置换分布。观测差异为 3.832，置换均值为 2.803，95\% 分位数为 3.519，置换检验 $p=0.0146$，说明两类玻璃的成分关联结构整体存在显著差异。

\begin{table}[H]
  \centering
  \caption{高钾与铅钡之间差异较显著的可靠成分相关关系}
  \label{tab:q4-diff}
  \begin{tabular}{lrrrr}
    \toprule
    成分对 & 高钾相关 & 铅钡相关 & 差值 & $q$ 值 \\
    \midrule
    $Fe_2O_3$--$CuO$ & 0.781 & -0.409 & 1.190 & 0.0182 \\
    $CaO$--$CuO$ & 0.622 & -0.409 & 1.030 & 0.0227 \\
    $Al_2O_3$--$PbO$ & 0.116 & -0.441 & 0.557 & 0.0394 \\
    $PbO$--$BaO$ & 0.587 & 0.196 & 0.391 & 0.0728 \\
    $K_2O$--$Al_2O_3$ & -0.220 & 0.341 & -0.562 & 0.0801 \\
    $SiO_2$--$Fe_2O_3$ & -0.665 & 0.099 & -0.765 & 0.0801 \\
    \bottomrule
  \end{tabular}
\end{table}

表 \ref{tab:q4-diff} 表明，差异主要来自 $Fe_2O_3$--$CuO$、$CaO$--$CuO$ 和 $Al_2O_3$--$PbO$。其中 $Fe_2O_3$--$CuO$ 在高钾中为 0.781，在铅钡中为 -0.409，方向相反；$CaO$--$CuO$ 也由高钾中的正相关转为铅钡中的负相关；$Al_2O_3$--$PbO$ 在铅钡中为 -0.441，体现铅相关成分与铝硅矿物杂质的替代关系。

\subsection{风化修正影响检查}

最后比较原始文物级成分和风化修正文物级成分下的相关矩阵距离。高钾原始与修正相关矩阵差异为 1.720，铅钡为 1.364，均小于两类之间的 3.832，因此第四问的类型差异大于风化修正扰动。第四问最终输出为：两类玻璃相关结构整体显著不同，高钾更突出铁铜、钙铜等关系，铅钡更突出钡铜、铝铅和铝钡等关系。

\section{模型评价与改进}

\subsection{模型优点}

本文模型的优点在于四问使用同一套数据口径和风化修正结果，避免了前后问题各自建模造成的结论断裂。CLR 变换处理了成分总和约束，对数比阈值保留了高钾与铅钡的化学解释，最近中心、聚类和置换检验则分别用于复核分类、寻找亚类和判断相关结构差异。整体上，模型不是追求复杂分类器，而是优先给出可解释、可复核的判别流程。

\subsection{模型缺点}

模型本身也存在局限。第一，风化前预测采用类型内平均风化效应，本质上是线性反向修正，默认同一类型内各成分的风化偏移方向相近，不能刻画成分迁移中的非线性、交互作用和局部饱和现象。第二，对数比阈值模型把高钾与铅钡的分界压缩到一个一维指标上，解释清楚但表达能力有限；若出现混合配方或过渡型样本，该模型可能只能给出硬分类，难以表达连续过渡程度。第三，聚类模型需要预先选择特征、距离和类数，亚类标签依赖这些建模选择，因而更适合解释内部结构，而不宜被理解为绝对类别。

相关结构比较也有方法边界。CLR 相关矩阵只能描述两两线性关联，不能直接说明因果关系，也不能区分“共同原料来源”“工艺配比约束”和“风化共同迁移”三类机制。Frobenius 差异能够概括两类相关矩阵的整体距离，但会把所有成分对差异压缩为一个标量，容易掩盖局部成分网络中的方向性信息。因此，第四问的相关结果更适合作为类别差异的结构证据，而不是单独推出工艺因果链。

\subsection{改进方向}

进一步改进可围绕模型形式展开：风化校正可由平均偏移模型扩展为成分迁移的非线性模型；大类判别可在对数比阈值之外增加概率输出，给边界样本提供置信度；亚类划分可采用模型化聚类或模糊聚类，使样本同时具有类别归属和不确定性；相关分析可引入偏相关网络或成分数据专用图模型，从整体相关矩阵进一步转向可解释的成分关系网络。

\begin{thebibliography}{99}
  \bibitem{aitchison1982} Aitchison J. The statistical analysis of compositional data[J]. Journal of the Royal Statistical Society: Series B, 1982, 44(2): 139--177.
  \bibitem{firth1993} Firth D. Bias reduction of maximum likelihood estimates[J]. Biometrika, 1993, 80(1): 27--38.
  \bibitem{bh1995} Benjamini Y, Hochberg Y. Controlling the false discovery rate: a practical and powerful approach to multiple testing[J]. Journal of the Royal Statistical Society: Series B, 1995, 57(1): 289--300.
  \bibitem{ward1963} Ward J H. Hierarchical grouping to optimize an objective function[J]. Journal of the American Statistical Association, 1963, 58(301): 236--244.
  \bibitem{macqueen1967} MacQueen J. Some methods for classification and analysis of multivariate observations[C]//Proceedings of the Fifth Berkeley Symposium on Mathematical Statistics and Probability. 1967: 281--297.
  \bibitem{openai2026} OpenAI. Codex based on GPT-5, AI assistant used for text drafting, code organization and consistency checking, 2026-08-05.
\end{thebibliography}

\CumcmMainMatterEnd
\appendix
\ctexset{section={name={附录,、},number=\Alph{section}}}

\section{支撑材料与代码文件说明}

\begin{longtable}{L{0.28\textwidth}L{0.14\textwidth}L{0.46\textwidth}}
  \caption{支撑材料文件清单}\label{tab:support-files}\\
  \toprule
  文件名 & 对应问题 & 用途说明 \\
  \midrule
  \endfirsthead
  \toprule
  文件名 & 对应问题 & 用途说明 \\
  \midrule
  \endhead
  \supportfile{source-workbook.xlsx}{全题}{题目附件数据的工作簿副本，包含已分类文物、采样点和未知文物。}
  \supportfile{cleaned-sampling-points.csv}{问题一}{有效采样点清洗结果，包含文物类别、外观信息、风化状态和 14 种成分。}
  \supportfile{artifact-state-compositions.csv}{问题一}{文物--状态分析单元，用于风化与无风化成分差异比较。}
  \supportfile{weathering-rates.csv}{问题一}{类型、纹饰、颜色的原始风化率。}
  \supportfile{firth-adjusted-probabilities.csv}{问题一}{多因素修正后的风化概率及区间。}
  \supportfile{component-difference-tests.csv}{问题一}{分类型成分差异置换检验与多重比较结果。}
  \supportfile{preweathering-predictions.csv}{问题一}{风化点风化前成分预测结果。}
  \supportfile{corrected-artifact-compositions.csv}{问题二}{风化修正后的文物级成分，是问题二至问题四的主要输入。}
  \supportfile{major-type-classification-summary.csv}{问题二}{高钾与铅钡大类分类模型性能。}
  \supportfile{subclass-assignments.csv}{问题二}{每件已知类别文物的亚类划分结果。}
  \supportfile{subclass-centers.csv}{问题二}{各亚类中心成分及命名。}
  \supportfile{unknown-final-identification.csv}{问题三}{未知文物最终类型和亚类鉴别结果。}
  \supportfile{unknown-type-sensitivity.csv}{问题三}{未知文物大类判别的伪计数敏感性结果。}
  \supportfile{type-conditional-weathering-classification.csv}{问题三}{风化未知样本的类型条件修正复核结果。}
  \supportfile{clr-correlation-high-k.csv}{问题四}{高钾玻璃 CLR 成分相关矩阵。}
  \supportfile{clr-correlation-lead-barium.csv}{问题四}{铅钡玻璃 CLR 成分相关矩阵。}
  \supportfile{differential-correlations.csv}{问题四}{两类玻璃相关差异的置换检验结果。}
  \supportfile{global-correlation-difference.csv}{问题四}{相关矩阵整体差异的置换检验结果。}
  \supportfile{AI 工具使用详情.pdf}{全题}{AI 工具使用详情，记录使用环节、采纳内容和人工核验情况。}
  \bottomrule
\end{longtable}

\begin{longtable}{L{0.24\textwidth}L{0.14\textwidth}L{0.50\textwidth}}
  \caption{可运行代码文件清单}\label{tab:code-files}\\
  \toprule
  文件名 & 对应问题 & 入口、输入与输出说明 \\
  \midrule
  \endfirsthead
  \toprule
  文件名 & 对应问题 & 入口、输入与输出说明 \\
  \midrule
  \endhead
  \supportfile{run_question1.py}{问题一}{读取原始附件，完成数据清洗、风化关联检验、成分差异检验和风化前预测，输出问题一结果表。}
  \supportfile{run_question2.py}{问题二}{读取问题一风化前预测和原始附件，建立大类分类规则、亚类聚类和敏感性分析，输出问题二结果表和图。}
  \supportfile{run_question3.py}{问题三}{读取表单 3 未知文物、问题二分类规则和问题一风化效应，输出未知文物鉴别结果和敏感性表。}
  \supportfile{run_question4.py}{问题四}{读取问题二文物级修正成分，计算 CLR 相关矩阵、差异检验和热图，输出问题四结果。}
  \bottomrule
\end{longtable}

完整源代码、数据表和 SVG 图见支撑材料目录。正文只列关键表格，完整长表用于复现和核查。

\end{document}
